\documentclass[a4]{article}

% Pacotes essenciais
\usepackage[utf8]{inputenc}
\usepackage[T1]{fontenc}
\usepackage[brazil]{babel}
\usepackage{geometry}
\usepackage{graphicx}
\usepackage{titling}
\usepackage{fancyhdr}
\usepackage{enumitem}
\usepackage{titlesec}
\usepackage{xcolor}
\usepackage{tcolorbox}
\usepackage{hyperref}
\usepackage{tocloft}
\usepackage{amsfonts}
\usepackage{tikz}
\usepackage{array}
\usepackage{everypage}
\usepackage{amsfonts}
\usepackage{amsmath}
\usepackage{xcolor}
\usepackage{circuitikz}
\usepackage{etoolbox} % Necessário para patching
\usepackage[backend=biber,style=abnt]{biblatex}
\addbibresource{ref.bib}


% Redefine o \subsection para adicionar uma linha horizontal
\preto{\subsection}{\vspace{1em}\noindent\textcolor{cyan!70!black}{\rule{\linewidth}{0.5pt}}\par\vspace{0.5em}}

\preto{\subsubsection}{\vspace{1em}\noindent\textcolor{cyan!70!black}{\rule{\linewidth}{0.5pt}}\par\vspace{0.5em}}

\definecolor{MidnightBlue}{RGB}{25,25,112} % ou a cor exata usada na sua barra
\definecolor{customblue}{RGB}{0, 0, 255}
\definecolor{customgreen}{RGB}{0,190, 255}
\definecolor{primary}{RGB}{20, 50, 90}
\definecolor{secondary}{RGB}{230, 240, 250}
\definecolor{correctgreen}{RGB}{34, 139, 34}
\definecolor{alertred}{RGB}{178, 34, 34}
% Margens
\geometry{a4paper, margin=2.5cm}

% Cabeçalho e rodapé
\pagestyle{fancy}
\fancyhf{}\rhead{\textit{\textcolor{customblue}{Lista1}}}
%\lhead{\textit{\textcolor{customblue}{Controle digital de aeronaves}}}%alterar
\fancyfoot[C]{\thepage}
\fancyhead[L]{\textit{\textcolor{customblue}{\nouppercase{\leftmark}}}} 

% Contador de página para ativar barras só a partir da 3ª página
\newcounter{mypage}

\AddEverypageHook{%
  \stepcounter{mypage}%
  \ifnum\themypage>1
    % Barra superior
    \begin{tikzpicture}[remember picture,overlay]
      \fill[cyan!15] (current page.north west) rectangle ([yshift=-1.8cm]current page.north east);
      \node[anchor=north, text=black, font=\bfseries\large] at ([yshift=-0.9cm]current page.north) 
      {};
    \end{tikzpicture}
    % Barra inferior
    \begin{tikzpicture}[remember picture,overlay]
      \fill[cyan!15] (current page.south west) rectangle ([yshift=2.0cm]current page.south east);
      \node[anchor=south, text=black, font=\bfseries\small] at ([yshift=0.9cm]current page.south) {\textit{\textcolor{customblue}{Universidade Federal de Santa Catarina }}};
    \end{tikzpicture}
  \fi
}
% Formato das seções
\titleformat{\section}{\bfseries\large}{\thesection.}{1em}{}

\usepackage{tcolorbox}
\tcbuselibrary{listings, breakable}

\newtcolorbox{MathBox}{
  colback=cyan!10,
  colframe=cyan!50!black,
  boxrule=0.5pt,
  arc=4pt,
  left=4pt,
  right=4pt,
  top=4pt,
  bottom=4pt
}

%Estilização das seções e subseções
\titleformat{\section}
  {\color{MidnightBlue}\normalfont\Large\bfseries}
  {\thesection}{1em}{}
\titleformat{\subsection}
  {\color{customblue}\normalfont\large\bfseries}
  {\thesubsection}{1em}{}
\titleformat{\subsubsection}
  {\color{customblue}\normalfont\large\bfseries}
  {\thesubsection}{1em}{}


\let\oldtextbf\textbf

\renewcommand{\textbf}[1]{\textcolor{customgreen}{\oldtextbf{#1}}}

% Estilização das listas
\setlist[itemize]{label=\textcolor{customblue}{\textbullet}}
\setlist[enumerate]{label=\textcolor{customblue}{\Roman*.}}

% Estilo de sumário
\cftsetindents{section}{0em}{2em}
\cftsetindents{subsection}{2em}{3em}
\cftsetindents{subsubsection}{5em}{4em}


% Capa
\begin{document}
\begin{titlepage}

    \centering
    \vspace*{3cm}
    {\Huge\bfseries Lista 2 - Instrumentação\\[1.5em]}

    {\large\textit{}}\\[7em]
    \includegraphics[width=0.3\textwidth]{ufsc}\\[1em]
    {\Large \textit{Instrumentação}}
    \vfill
    %\today
\end{titlepage}

\newpage

\begin{center}
    \Large\textbf{Contributors}\\[2em]
    
    \begin{itemize}
        \item \textbf{Marcos Antonio Tomé Oliveira} \\\textit{Graduando em engenharia mecatrônica}\\[0.5em]
       
    \end{itemize}

    %\textbf{Revisado em:} 12 de dezembro de 2023\\
    %\textbf{Organização:} RoboCore\\[6em]
    %\textbf{Link:} \href{https://events.robocore.net}{https://events.robocore.net}
\end{center}
\newpage
\begin{enumerate}[label=\textbf{\arabic*)}, leftmargin=*]

    \item 
    \begin{itemize}
        \item \textbf{transdutor:} Transforma um tipo de energia em outro, podem ser encontrados em estágio de entrada e saída.
        \item \textbf{sensor:} capta uma grandeza física.
        \item \textbf{atuador:} Transforma um sinal elétrico em outra grandeza física.
    \end{itemize}

    \item Um transdutor ativo possui alimentação externa própria, enquanto um transdutor passivo não necessita de energia adicional.

    \item Um sinal analógico possui um sinal contínuo, ou seja, sua curva não possui bicos, já o sinal digital varia entre zeros e uns.

    \item Para a saída em corrente, o instrumento/sensor gera um sinal em forma de corrente elétrica; para saída em tensão, o sensor gera saída em tensão.

    \item Em um instrumento caracterizado como simples há apenas uma conversão no sinal de entrada, no instrumento composto há mais de uma conversão no sinal de entrada.

    \item No método de deflexão a grandeza provoca uma resposta direta (deflexão) no instrumento, e essa resposta é usada para medir a grandeza. No método de detecção de erro nulo, a grandeza é comparada a um sinal de referência ajustável, até que o erro seja nulo.

    \item Os erros aleatórios interferem em cada medida de maneira aleatória, enquanto erros sistemáticos acontecem de maneira semelhante.
    \item \textbf{Análise do alvo (Exatidão / Precisão):}
    \begin{enumerate}[label=\alph*)]
        \item Alta exatidão e alta precisão.
        \item Baixa exatidão e alta precisão.
        \item Alta exatidão e baixa precisão.
        \item Baixa exatidão e baixa precisão.
    \end{enumerate}

    \item \textbf{Algarismos significativos:}
    \begin{enumerate}[label=\alph*)]
        \item $4{,}35 \rightarrow 3$ algarismos significativos.
        \item $0{,}01 \rightarrow 1$ algarismo significativo.
        \item $0{,}0000052 \rightarrow 2$ algarismos significativos.
        \item $400{,}510 \rightarrow 6$ algarismos significativos.
        \item $1000{,}1 \rightarrow 5$ algarismos significativos.
        \item $0{,}00000000000001 \rightarrow 1$ algarismo significativo.
    \end{enumerate}

    \item \textbf{Arredondamento:}
    \begin{enumerate}[label=\alph*)]
        \item $2{,}2\bar{4}5 \rightarrow 2{,}25$
        \item $6\bar{9},87 \rightarrow 70$
        \item $123,\bar{1}2 \rightarrow 123{,}1$
    \end{enumerate}

    \item \textbf{Operações com algarismos significativos:}
    \begin{enumerate}[label=\alph*)]
        \item $12{,}45 \times 7{,}2 = 89{,}64 \approx 90$
        \item $5{,}2 + 7{,}25 = 12{,}45 \approx 12{,}5$
        \item $351{,}12 - 136 = 215{,}12 \approx 215$
        \item $134{,}15 / 2{,}15 = 6{,}2395 \approx 6{,}24$
        \item $(12{,}10 + 5{,}1) \times (6{,}75 - 4{,}89) = 17{,}2 \times 1{,}86 = 31{,}992 \approx 32$
    \end{enumerate}

    \item
    \begin{enumerate}[label=\alph*)]
        \item \textbf{Médias aritméticas:}
        \[
        \bar{P} = \frac{24{,}95 + 25{,}34 + 25{,}14 + 24{,}86 + 25{,}43}{5} = 25{,}14 \approx 25
        \]
        \[
        \bar{i} = \frac{2{,}95 + 3{,}05 + 3{,}00 + 2{,}98 + 3{,}10}{5} = 3{,}02 \approx 3
        \]

        \item \textbf{Desvio padrão da média:}
        \[
        \sigma_P = 0{,}24419 \implies \sigma_{\bar{P}} = \frac{0{,}24419}{\sqrt{5}} = 0{,}1092 = \Delta P
        \]
        \[
        \sigma_I = 0{,}05941 \implies \sigma_{\bar{I}} = \frac{0{,}05941}{\sqrt{5}} = 0{,}02657 = \Delta I
        \]

        \item \textbf{Cálculo da Tensão ($V = P / I$) e Incerteza Propagada:}
        \[
        \bar{V} = \frac{\bar{P}}{\bar{I}} = \frac{25}{3} = 8{,}33\,\text{V} \approx 8\,\text{V}
        \]
        
        Derivadas parciais:
        \[
        \frac{\partial V}{\partial P} = \frac{1}{I} = \frac{1}{3}
        \]
        \[
        \frac{\partial V}{\partial I} = -P \cdot I^{-2} = -25 \cdot 3^{-2} = -2{,}778
        \]

        Propagação da incerteza:
        \[
        \Delta V = \left| \frac{\partial V}{\partial P} \right| \Delta P + \left| \frac{\partial V}{\partial I} \right| \Delta I
        \]
        \[
        \Delta V = \frac{1}{3} \cdot 0{,}11 + 2{,}77 \cdot 0{,}027 = 0{,}112
        \]

        Cálculo do valor médio direto de $V$:
        \[
        \bar{V} = \frac{8{,}46 + 8{,}31 + 8{,}38 + 8{,}34 + 8{,}20}{5} = 8{,}34 \approx 8
        \]
        \[
        \sigma_{\bar{V}} = 0{,}08541
        \]

        \item \textbf{Resultados no formato padrão:}
        \begin{align*}
        V &= 8{,}34 \pm 0{,}09 \\
        i &= 3{,}02 \pm 0{,}05 \\
        P &= 25{,}14 \pm 0{,}01
        \end{align*}
    \end{enumerate}

\end{enumerate}

% \section*{Questão 2.29}


\subsection*{Modelo e relações úteis}

Para um Amp-Op de um único polo podemos escrever a resposta em malha aberta como
\[
A(j\omega)=\frac{A_0}{1+j\frac{\omega}{\omega_p}},
\]
onde $\omega_p=2\pi f_p$ e a pulsação do polo. A frequência de ganho unitário $f_t$ (gain-bandwidth) relaciona-se com $A_0$ e $f_p$ por
\[
f_t = A_0 f_p \quad\Longrightarrow\quad f_p=\frac{f_t}{A_0}.
\]

No esquema inversor ideal com ganho em malha fechada $A_{CL}=-\dfrac{1}{\beta}$ temos $\beta = 1/|A_{CL}|$. Para um único polo, a função de transferência em malha fechada pode ser aproximada por uma resposta de primeira ordem:
\[
H(j\omega)\approx \frac{A_{CL}}{1 + j\frac{\omega}{\omega_B}},
\]
onde a largura de banda (pulsacao) em malha fechada $\omega_B$ corresponde a frequencia de −3 db de malha fechada e, para polo unico,
\[
f_B \equiv f_{3\text{dB}} \approx \frac{f_t}{|A_{CL}|}.
\]
(Essa fórmula e a conhecida relacao \emph{GBW / ganho = largura de banda} valida para amplificadores de um polo.)

\section*{Cálculos numericos}

\paragraph{Dados:}
\[
A_{CL} = -20,\qquad |A_{CL}|=20,\qquad A_0=10^5,\qquad f_t=10^6\ \mathrm{Hz}.
\]

\subsection*{(a) Frequência de −3 dB em malha fechada}
\[
f_{3\text{dB}} = \frac{f_t}{|A_{CL}|} = \frac{10^{6}}{20} = 5\times 10^{4}\ \mathrm{Hz}=50\ \text{kHz}.
\]

\subsection*{(b) Ganho em $0{,}1\,f_t$ e em $10\,f_t$}

A expressão de magnitude (para o modelo de 1º ordem) e
\[
|H(f)| = \frac{|A_{CL}|}{\sqrt{1+\left(\dfrac{f}{f_{3\text{dB}}}\right)^2}}.
\]

\paragraph*{Para } $f = 0{,}1\,f_t = 0{,}1\times 10^{6}=10^{5}\ \mathrm{Hz}$:
\[
\frac{f}{f_{3\text{dB}}}=\frac{10^{5}}{5\times10^{4}}=2,
\]
\[
|H|=\frac{20}{\sqrt{1+2^{2}}}=\frac{20}{\sqrt{5}}\approx 8{,}944\ \mathrm{V/V}.
\]
Em dB:
\[
20\log_{10}(8{,}944)\approx 19{,}03\ \mathrm{dB}.
\]

\paragraph*{Para } $f = 10\,f_t = 10\times 10^{6}=10^{7}\ \mathrm{Hz}$:
\[
\frac{f}{f_{3\text{dB}}}=\frac{10^{7}}{5\times10^{4}}=200,
\]
\[
|H|\approx\frac{20}{\sqrt{1+200^{2}}}\approx\frac{20}{200}=0{,}1\ \mathrm{V/V}.
\]
Em dB:
\[
20\log_{10}(0{,}1) = -20\ \mathrm{dB}.
\]
\subsection*{Vizualização dos resultados}

\begin{figure}[h!]
    \centering
    \includegraphics[width=0.5\linewidth]{2.29.png}
    \caption{Magnitude do Amplificador}
    \label{fig:placeholder}
\end{figure}
\section*{Observações finais}
\begin{itemize}
  \item O ganho DC em dB e $20\log_{10}(20)\approx 26{,}02\ \mathrm{dB}$. Em $f_{3\text{dB}}$ a magnitude e $20/\sqrt{2}\approx 14{,}14$ (≈23,02 dB), ou seja, 3 dB abaixo do ganho DC — isto confirma a interpretacao de $f_{3\text{dB}}$.
  \item Para $f\gg f_{3\text{dB}}$ a resposta decai aproximadamente a $-20\ \mathrm{dB/dec}$ (comportamento de um polo único), e o valor em $f_t$ tende a $\approx 1$ (0 dB), já que $f_t/ f_{3\text{dB}} = |A_{CL}|$.
\end{itemize}
\section*{Questão 2.34}

O amplificador operacional possui \textit{slew rate}
\[
SR = 10\ \text{V}/\mu\text{s} = 10\times 10^{6}\ \text{V/s}
\]
e a onda triangular possui amplitude pico a pico
\[
V_{pp} = 20\ \text{V}.
\]

Para uma onda triangular, a máxima inclinação ocorre nos trechos lineares e é dada por
\[
\left|\frac{dv}{dt}\right|_{\max} = 2 V_{pp} f.
\]
Impondo a condição para ausência de distorção:
\[
2V_{pp} f_{\max} = SR
\]
\[
f_{\max} = \frac{SR}{2V_{pp}}
= \frac{10\times 10^{6}}{40}
= 2.5\times 10^{5}\ \text{Hz}
= 250\ \text{kHz}.
\]

\subsection*{Senoide na mesma frequência}

Para uma senoide:
\[
v(t)=V_p\sin(2\pi f t)
\]
\[
\left|\frac{dv}{dt}\right|_{\max}=V_p 2\pi f.
\]
O limite imposto pelo \textit{slew rate} é
\[
V_p 2\pi f = SR,
\]
portanto:
\[
V_{p,\max} = \frac{SR}{2\pi f}
= \frac{10\times 10^{6}}{2\pi (250\times10^{3})}
= \frac{20}{\pi}
\approx 6.366\ \text{V}.
\]

A amplitude pico a pico é:
\[
V_{pp,\max}=2V_p\approx 12.73\ \text{V}.
\]

---

\section*{Questão 2.42}

O amplificador operacional possui tensão de offset:
\[
v_{os}=4\ \text{mV}.
\]
Como a entrada não inversora está aterrada, então:
\[
V_+ = 0,\qquad V_- = -v_{os}.
\]

\subsection*{Caso 1: Entrada ligada diretamente (sem capacitor)}

Aplicamos KCL no nó $V_-$:
\[
\frac{V_- - 0}{1\,\text{M}\Omega}
+ \frac{V_- - V_A}{1\,\text{M}\Omega} = 0.
\]
Assim:
\[
2V_- - V_A = 0
\]
\[
V_A = 2V_-.
\]

Agora aplicamos KCL no nó $V_A$:
\[
\frac{V_A - V_-}{1\,\text{M}\Omega}
+ \frac{V_A - 0}{1\,\text{k}\Omega}
+ \frac{V_A - V_{out}}{1\,\text{M}\Omega}=0.
\]

Substituindo $V_A=2V_-$:
\[
\frac{2V_- - V_-}{1\,\text{M}\Omega}
+ \frac{2V_-}{1\,\text{k}\Omega}
+ \frac{2V_- - V_{out}}{1\,\text{M}\Omega}=0.
\]

Simplificando:
\[
\frac{V_-}{1\,\text{M}\Omega}+\frac{2V_-}{1\,\text{k}\Omega}
+ \frac{2V_- - V_{out}}{1\,\text{M}\Omega}=0.
\]

Agrupando termos e resolvendo:
\[
V_{out}\approx -8.01\ \text{V}.
\]

---

\subsection*{Caso 2: Entrada acoplada em CA (capacitor $\Rightarrow$ entrada DC aberta)}

No regime DC, a entrada fica aberta.

O nó $V_-$ conecta apenas aos resistores para $V_A$:
\[
V_- = V_A = -v_{os}.
\]

Aplicando KCL no nó $V_A$:
\[
\frac{V_A - V_-}{1\,\text{M}\Omega}
+\frac{V_A - V_{out}}{1\,\text{M}\Omega}
+\frac{V_A - 0}{1\,\text{k}\Omega}=0,
\]
mas como $V_A - V_- = 0$, sobra:
\[
\frac{V_A - V_{out}}{1\,\text{M}\Omega}
+\frac{V_A}{1\,\text{k}\Omega}=0.
\]

Multiplicando por $1\,\text{M}\Omega$:
\[
V_A - V_{out} + 1000\,V_A = 0.
\]

\[
V_{out} = 1001\,V_A = -1001\,v_{os}.
\]

\[
V_{out}\approx -4.004\ \text{V}.
\]

---

\subsection*{Caso 3: Resistor de 1 k$\Omega$ acoplado em CA (aberto no DC)}

Agora o nó $A$ não tem caminho DC para terra.

KCL no nó $V_-$:
\[
\frac{V_- - 0}{1\,\text{M}\Omega}
+ \frac{V_- - V_A}{1\,\text{M}\Omega}=0
\]
\[
2V_- - V_A = 0
\]
\[
V_A = 2V_-.
\]

KCL no nó $A$ (sem resistor para terra):
\[
\frac{V_A - V_-}{1\,\text{M}\Omega}
+ \frac{V_A - V_{out}}{1\,\text{M}\Omega}=0.
\]

Substituindo $V_A = 2V_-$:
\[
\frac{2V_- - V_-}{1\,\text{M}\Omega}
+ \frac{2V_- - V_{out}}{1\,\text{M}\Omega}=0
\]

\[
\frac{V_-}{1\,\text{M}\Omega}
+ \frac{2V_- - V_{out}}{1\,\text{M}\Omega}=0.
\]

Somando:
\[
3V_- - V_{out} = 0
\]
\[
V_{out}=3V_- = -3v_{os}.
\]

\[
V_{out}=-12\ \text{mV}.
\]

---

\section*{Questão 12.1}

A função de transferência é:
\[
T(j\omega)=\frac{\omega_0}{\sqrt{\omega^2+\omega_0^2}}
=\frac{1}{\sqrt{1+(\omega/\omega_0)^2}}.
\]

Vamos calcular os valores de $|T|$ e de $20\log_{10}|T|$ para:
\[
\omega/\omega_0 = 0,\ 0.5,\ 1,\ 10,\ 100.
\]

\subsection*{Cálculo ponto a ponto}

\[
|T(0)| = 1
\]

\[
|T(0.5\omega_0)| = \frac{1}{\sqrt{1+0.5^2}}
= 0.8944
\]

\[
|T(\omega_0)| = \frac{1}{\sqrt{2}}
= 0.7071
\]

\[
|T(10\omega_0)| = \frac{1}{\sqrt{101}}
= 0.09950
\]

\[
|T(100\omega_0)| = \frac{1}{\sqrt{10001}}
= 0.009999
\]

Os valores em decibéis:

\[
20\log_{10}|T(0)| = 0\ \text{dB}
\]
\[
20\log_{10}(0.8944) = -0.969\ \text{dB}
\]
\[
20\log_{10}(0.7071) = -3.010\ \text{dB}
\]
\[
20\log_{10}(0.09950) = -20.04\ \text{dB}
\]
\[
20\log_{10}(0.009999) = -40.00\ \text{dB}
\]

\subsection*{Tabela resumo}

\begin{table}[h!]
\centering
\begin{tabular}{c|c|c}
\hline
$\omega/\omega_0$ & $|T(j\omega)|$ & $20\log_{10}|T(j\omega)|$ (dB) \\
\hline
0   & 1.0000   & 0.00 \\
0.5 & 0.8944   & -0.97 \\
1   & 0.7071   & -3.01 \\
10  & 0.09950  & -20.04 \\
100 & 0.009999 & -40.00 \\
\hline
\end{tabular}
\end{table}


\end{document}

